\subsubsection{Subproceso 2: Selección y Registro de Cursos}



\textbf{Inicio:} El subproceso se activa cuando el estudiante accede a la Intranet DIRCE dentro de su franja horaria habilitada y el sistema confirma su estado de ``Habilitado'' (salida del Subproceso 1). En ese momento la plataforma despliega la oferta académica personalizada según su Plan de Estudios y su historial de prerrequisitos aprobados.

\textbf{Incluye:}
\begin{itemize}
    \item Visualización de la oferta de asignaturas, secciones, horarios y vacantes disponibles para el ciclo en curso.
    \item Selección interactiva de asignaturas y secciones por parte del estudiante.
    \item Validación automática en tiempo real del límite de créditos según promedio ponderado y de la inexistencia de cruces de horario prohibidos.
    \item Registro y envío de la propuesta de matrícula mediante confirmación final en el portal.
    \item Generación de la reserva temporal de cupos en el SIGA-ORCE.
\end{itemize}

\textbf{Fin:} El subproceso concluye cuando el sistema registra la Pre-Ficha de Matrícula con las asignaturas y secciones seleccionadas y emite el comprobante provisional al estudiante, transfiriendo la propuesta al Subproceso 3.A para su auditoría.

\textbf{Fuera del alcance:}
\begin{itemize}
    \item La verificación de compromisos financieros (autoseguro, deudas) y la asignación del turno, que corresponden al Subproceso 1.
    \item La auditoría de prerrequisitos y reglas académicas de fondo, que se realiza en el Subproceso 3.A.
    \item La confirmación transaccional definitiva y emisión de la Boleta Oficial, propias del Subproceso 3.B.
    \item La resolución de solicitudes de rectificación por causales de vacantes agotadas o cruces, gestionadas en el Subproceso 3.A.
\end{itemize}




